\documentclass[12pt,a4paper]{article}
\usepackage[UTF8]{ctex}
\usepackage{amsmath,amssymb}
\usepackage{graphicx}
\usepackage{booktabs}
\usepackage{geometry}
\usepackage{hyperref}
\usepackage{caption}
\usepackage{subcaption}
\usepackage{multirow}
\usepackage{setspace}
\geometry{margin=2.5cm}
\onehalfspacing

\title{\textbf{融合上游水文信息的注意力 LSTM 河流流量多步预报研究}\\[0.5em]
\large ——以波托马克河流域为例}
\author{作者姓名$^{1}$ \quad 通讯作者$^{1,*}$\\
\small $^{1}$某某大学 水利与环境学院，城市 邮编\\
\small $^{*}$通讯作者：email@university.edu}
\date{}

\begin{document}
\maketitle

\begin{abstract}
准确的中短期河流流量预报对防洪减灾与水资源调度至关重要。本研究以美国波托马克河流域为案例，利用美国地质调查局（USGS）2000--2023年日尺度流量实测数据及 Open-Meteo 气象再分析资料，构建了融合上游支流流量的 LSTM-Attention 深度学习预报模型。模型输入包含目标站、两处上游水文站的流量与降雨序列及气温，输出未来 1、3、7 天流量预报。与持续性预报、ARIMA、XGBoost 及标准 LSTM 对比表明：在 1 天预报尺度，LSTM-Attention 的 Nash--Sutcliffe 效率系数（NSE）达 0.930，Kling--Gupta 效率系数（KGE）为 0.884，RMSE 为 2422 cfs，优于全部基线模型；XGBoost（NSE=0.916）与 LSTM（NSE=0.912）亦表现优异。在 3 天尺度，LSTM（NSE=0.543）略优于其他方法；7 天尺度所有模型 NSE 均低于 0.11，表明周尺度洪枯转换预报仍是难点。消融实验证实，引入上游水文站信息使 1 天预报 NSE 从 0.889 提升至 0.923。研究为数据驱动洪水预报提供了可复现的基准框架，并揭示了多步预报中误差累积与物理信息融合的关键作用。

\textbf{关键词：}流量预报；LSTM；注意力机制；波托马克河；深度学习；水文时序
\end{abstract}

\begin{abstract}
\noindent\textbf{Abstract:} Accurate short- to medium-term streamflow forecasting is critical for flood mitigation and water resources management. This study develops an LSTM-Attention model that integrates upstream tributary streamflow for the Potomac River basin using USGS daily discharge records (2000--2023) and Open-Meteo meteorological reanalysis. The model ingests multi-station flow, precipitation, and temperature series to produce 1-, 3-, and 7-day ahead forecasts at Washington, D.C. (USGS-01646500). Compared with persistence, ARIMA, XGBoost, and standard LSTM, LSTM-Attention achieves the best 1-day performance (NSE = 0.930, KGE = 0.884, RMSE = 2422 cfs). At the 3-day horizon, LSTM performs best (NSE = 0.543). All models degrade substantially at 7 days (NSE $<$ 0.11), highlighting the difficulty of weekly forecasting. Ablation experiments show that upstream gauge information improves 1-day NSE from 0.889 to 0.923. A fully reproducible pipeline is provided for benchmark hydrological machine learning research.

\noindent\textbf{Keywords:} streamflow forecasting; LSTM; attention mechanism; Potomac River; deep learning; hydrological time series
\end{abstract}

\section{引言}

全球气候变化与人类活动加剧了极端水文事件的频率与强度，精准的流量预报已成为水文学研究的核心问题之一\cite{beven2011}。传统水文模型（如概念性新安江模型、HEC-HMS）基于物理机制与参数率定，在资料匮乏地区往往面临结构误差与参数不确定性\cite{kirchner2006}。近年来，深度学习尤其是长短期记忆网络（LSTM）在捕捉非线性时序依赖方面展现出显著优势\cite{kratzert2019,frame2022}。

然而，现有研究仍存在三方面不足：（1）多数工作仅利用单站历史流量，未充分利用流域内上下游水力联系；（2）注意力机制虽可增强模型可解释性，但其在多步预报中的稳定性缺乏系统评估；（3）与经典统计模型及树模型的公平对比尚不充分。针对上述问题，本研究以波托马克河下游控制站为预报目标，融合 Shenandoah 河与 Harpers Ferry 两处上游站流量，构建 LSTM-Attention 模型，系统评估 1--7 天预报性能，并通过消融实验量化上游信息贡献。

\section{研究区与数据}

\subsection{研究流域}

波托马克河是美国东部大西洋沿岸最大河流之一，流域面积约 14,670 km$^2$，流经马里兰、弗吉尼亚、西弗吉尼亚等州，在华盛顿特区汇入切萨皮克湾。本研究选取三处 USGS 水文站构成由上游至下游的监测网络（图~\ref{fig:basin}）：

\begin{itemize}
  \item \textbf{01636500}：Shenandoah River at Millville, WV（支流）
  \item \textbf{01638480}：Potomac River at Harpers Ferry, WV（中游）
  \item \textbf{01646500}：Potomac River at Washington, DC（下游控制站，预报目标）
\end{itemize}

\subsection{数据来源与预处理}

\textbf{流量数据}来自美国地质调查局国家水质信息系统（USGS NWIS）日平均流量接口（参数代码 00060，单位 cfs），时间跨度 2000-01-01 至 2023-12-31，共 8764 天。

\textbf{气象数据}来自 Open-Meteo Historical Weather API 的日降水（mm）与 2 m 气温（$^\circ$C），按各站经纬度提取。

\textbf{数据划分}采用严格时间顺序切分，避免信息泄露：
\begin{itemize}
  \item 训练集：2000--2018（6940 天）
  \item 验证集：2019--2020（731 天）
  \item 测试集：2021--2023（1093 天）
\end{itemize}

缺失值剔除后，三站同步观测记录 8764 天。所有特征经 StandardScaler 标准化，标准化参数仅由训练集拟合。

\section{方法}

\subsection{问题定义}

给定长度为 $L=30$ 天的多变量输入序列 $\mathbf{X}_t = \{x_{t-L+1}, \ldots, x_t\}$，其中 $x_t \in \mathbb{R}^{d}$ 包含目标站及上游站的流量、降雨和气温，预测未来 $h \in \{1,3,7\}$ 天目标站流量：
\begin{equation}
  \hat{Q}_{t+h} = f_\theta(\mathbf{X}_t)
\end{equation}

\subsection{模型结构}

\textbf{LSTM}：双层 LSTM（隐藏维度 64，dropout 0.2），取末时刻隐状态经全连接层输出。

\textbf{LSTM-Attention}：在 LSTM 输出序列 $\{h_1,\ldots,h_L\}$ 上施加时间注意力：
\begin{equation}
  \alpha_i = \frac{\exp(w^\top h_i)}{\sum_{j=1}^{L}\exp(w^\top h_j)}, \quad
  c = \sum_{i=1}^{L}\alpha_i h_i, \quad \hat{Q}_{t+h} = W c + b
\end{equation}
注意力权重 $\alpha_i$ 可揭示历史关键时段。

\textbf{基线模型}：
\begin{itemize}
  \item \textit{Persistence}：$\hat{Q}_{t+h} = Q_t$
  \item \textit{ARIMA(1,0,1)}：每 60 天重率定，滚动多步预报
  \item \textit{XGBoost}：将 $L \times d$ 维序列展平为特征向量
  \item \textit{LSTM}：不含注意力机制
\end{itemize}

\subsection{训练与评估}

损失函数为 MSE，优化器 Adam（$lr=10^{-3}$），批大小 64，最大 80 epoch，以验证集损失早停。评价指标采用水文领域通用的 NSE、KGE、RMSE、MAE 和 PBIAS\cite{nash1970,gupta2009}：

\begin{equation}
  \text{NSE} = 1 - \frac{\sum(Q_{\text{obs}}-Q_{\text{sim}})^2}{\sum(Q_{\text{obs}}-\bar{Q}_{\text{obs}})^2}
\end{equation}

NSE $>$ 0.5 通常视为可接受，$>$ 0.75 为良好\cite{moriasi2007}。

\subsection{消融实验}

对比两种输入配置：（A）融合全部上游站流量与降雨及 7/30 天累积降雨；（B）仅目标站本地流量、降雨与气温。均采用 LSTM-Attention 架构。

\section{结果}

\subsection{多模型多步预报对比}

表~\ref{tab:metrics} 汇总测试集（2021--2023）结果。

\begin{table}[htbp]
\centering
\caption{各模型在不同预报尺度下的性能对比（测试集）}
\label{tab:metrics}
\begin{tabular}{llccccc}
\toprule
预报尺度 & 模型 & NSE & KGE & RMSE (cfs) & MAE (cfs) & PBIAS (\%) \\
\midrule
\multirow{5}{*}{1 天} 
  & Persistence & 0.787 & 0.893 & 4226 & 1415 & 0.1 \\
  & ARIMA       & $-$0.808 & 0.265 & 12305 & 5986 & 20.3 \\
  & XGBoost     & 0.916 & 0.945 & 2650 & 1035 & 3.5 \\
  & LSTM        & 0.912 & 0.915 & 2716 & 988 & $-$0.9 \\
  & LSTM-Attention & \textbf{0.930} & 0.884 & \textbf{2422} & 1014 & $-$5.2 \\
\midrule
\multirow{5}{*}{3 天}
  & Persistence & 0.164 & 0.584 & 8366 & 3296 & 0.5 \\
  & ARIMA       & $-$0.423 & 0.313 & 10914 & 5893 & 26.4 \\
  & XGBoost     & 0.472 & 0.632 & 6647 & 3349 & 15.2 \\
  & LSTM        & \textbf{0.543} & \textbf{0.654} & \textbf{6183} & \textbf{2892} & 8.3 \\
  & LSTM-Attention & 0.459 & 0.612 & 6729 & 3226 & 1.1 \\
\midrule
\multirow{5}{*}{7 天}
  & Persistence & $-$0.393 & 0.342 & 10802 & 4879 & 2.9 \\
  & ARIMA       & $-$0.087 & 0.240 & 9539 & 5996 & 30.3 \\
  & XGBoost     & \textbf{0.108} & \textbf{0.383} & \textbf{8641} & 5251 & 27.0 \\
  & LSTM        & 0.096 & 0.366 & 8699 & 5390 & 26.1 \\
  & LSTM-Attention & $-$0.058 & 0.280 & 9412 & 6327 & 44.7 \\
\bottomrule
\end{tabular}
\end{table}

\textbf{1 天预报}：LSTM-Attention 取得最优 NSE（0.930），较持续性预报（0.787）提升 18.2\%，RMSE 降低 42.7\%。XGBoost 与 LSTM 亦达到 NSE $>$ 0.91，表明机器学习模型能有效学习降雨--径流非线性响应。ARIMA 在未引入降雨协变量时表现极差（NSE=$-$0.808），说明单变量统计模型难以适应洪水期剧变过程。

\textbf{3 天预报}：所有模型性能显著下降，LSTM（NSE=0.543）略优。LSTM-Attention 在此尺度未展现优势，可能与注意力权重在长预报步长下过拟合有关。

\textbf{7 天预报}：NSE 均趋近于零，最佳 XGBoost 仅 0.108。周尺度预报受气象预报不确定性、土壤湿度记忆效应等累积误差制约，需引入集合预报或物理约束。

图~\ref{fig:hydrograph} 展示了 7 天尺度测试期流量过程线对比（见 \texttt{figures/hydrograph\_7d.png}）。

\subsection{消融实验：上游信息贡献}

\begin{table}[htbp]
\centering
\caption{上游水文站信息消融实验（LSTM-Attention）}
\label{tab:ablation}
\begin{tabular}{lccc}
\toprule
输入配置 & 1 天 NSE & 3 天 NSE & 7 天 NSE \\
\midrule
融合上游站（完整） & \textbf{0.923} & \textbf{0.343} & \textbf{0.036} \\
仅本地站         & 0.889 & 0.337 & $-$0.024 \\
\bottomrule
\end{tabular}
\end{table}

引入 Shenandoah 与 Harpers Ferry 上游流量使 1 天 NSE 提升 0.034（3.8\%），RMSE 从 3044 降至 2543 cfs。上游信息通过水力传播时差为下游预报提供先导信号，验证了流域空间关联对短期预报的价值。

\section{讨论}

\subsection{模型适用性分析}

本研究 1 天预报结果（NSE $>$ 0.91）达到国际水文预报竞赛（如 CAMELS）中深度学习模型的先进水平\cite{kratzert2019}。LSTM-Attention 在短尺度优于标准 LSTM，说明注意力机制有助于识别洪水来临前的关键降雨--流量响应时段。然而，注意力并未在所有尺度上稳定占优，提示未来可探索 Temporal Fusion Transformer 等更先进的时序架构\cite{lim2021}。

\subsection{ARIMA 基线的局限性}

ARIMA 表现不佳主要归因于：（1）未纳入降雨等外生变量（未来可采用 ARIMAX/SARIMAX）；（2）波托马克河洪水过程高度非线性，线性自回归结构难以刻画；（3）周期性重率定策略尚不足以适应丰枯转换。该结果也印证了在水文非平稳条件下单纯统计模型的局限\cite{milner2020}。

\subsection{业务化展望}

本研究所用 USGS 与 Open-Meteo 数据均为公开 API，代码全流程可复现（见项目仓库）。业务部署需考虑：（1）实时气象预报替代观测降雨；（2）模型在线更新与漂移监测；（3）不确定性量化（如 MC Dropout 或分位数回归）以支撑风险决策。

\section{结论}

本研究构建了融合多站上游水文信息的 LSTM-Attention 流量预报框架，主要结论如下：

\begin{enumerate}
  \item 在波托马克河 Washington 站 1 天预报中，LSTM-Attention 取得 NSE=0.930，显著优于持续性预报与传统 ARIMA，与 XGBoost、LSTM 同属高性能梯队。
  \item 预报尺度延长至 3--7 天时，所有模型性能急剧退化，揭示中短期洪水预报的核心挑战在于误差累积。
  \item 消融实验证实上游支流流量对 1 天预报具有显著增益，支持"流域空间信息 + 深度学习"的技术路线。
  \item 提供基于 USGS 公开数据的可复现基准，可为后续图神经网络流域建模与物理约束神经网络研究奠定基础。
\end{enumerate}

\section*{数据与代码可用性声明}

流量数据：\url{https://waterservices.usgs.gov/}；气象数据：\url{https://open-meteo.com/}。完整代码与结果见项目目录 \texttt{hydro-ml-paper/}。

\begin{thebibliography}{99}
\bibitem{beven2011} Beven K. Deep learning, hydrological processes and the unity of hydrology[J]. Hydrology Process, 2020, 34(16): 3382-3383.
\bibitem{kirchner2006} Kirchner J W. Getting the right answers for the right reasons[J]. Hydrology Process, 2006, 20(18): 4065-4068.
\bibitem{kratzert2019} Kratzert F, et al. Towards learning universal, regional, and local hydrological behaviors via ML applied to large-sample dataset[J]. Hydrology Earth System Sciences, 2019, 23(12): 5089-5110.
\bibitem{frame2022} Frame J M, et al. Deep learning rainfall-runoff predictions of extreme events[J]. Hydrology Earth System Sciences, 2022, 26(13): 3377-3392.
\bibitem{nash1970} Nash J E, Sutcliffe J V. River flow forecasting through conceptual models[J]. Journal of Hydrology, 1970, 10(3): 282-290.
\bibitem{gupta2009} Gupta H V, et al. Decomposition of the mean squared error and NSE criteria[J]. Water Resources Research, 2009, 45(3): W00B07.
\bibitem{moriasi2007} Moriasi D N, et al. Model evaluation guidelines for systematic quantification of accuracy[J]. Transactions of the ASABE, 2007, 50(3): 885-900.
\bibitem{lim2021} Lim B, et al. Temporal Fusion Transformers for interpretable multi-horizon time series forecasting[J]. International Journal of Forecasting, 2021, 37(4): 1748-1764.
\bibitem{milner2020} Milner A M, et al. Nonstationary streamflow response in a changing watershed[J]. Environmental Research Letters, 2020, 15(11): 114032.
\end{thebibliography}

\end{document}
